%! AP Statistics — Unit 4 Cover Sheet
\documentclass[10pt]{article}
\usepackage{apstats-article}
\usepackage{apstats-boxes}

%=============================================================================
\begin{document}

% ── Full-bleed navy banner ────────────────────────────────────────────────────
\begin{tikzpicture}[remember picture, overlay]
  \fill[navy] (current page.north west)
    rectangle ([yshift=-1.4in]current page.north east);
\end{tikzpicture}
\vspace{-1.3in}

\begin{center}
  {\color{white}\bfseries\fontsize{28}{34}\selectfont AP Statistics}\\[8pt]
  {\color{goldacc}\bfseries\Large Shepherd \quad 2026--2027}\\[14pt]
  {\color{white}\bfseries\LARGE Unit 4: Probability, Random Variables, and Probability Distributions}\\[4pt]
\end{center}

\vspace{0.30in}

% ── Unit overview ─────────────────────────────────────────────────────────────
\begin{tcolorbox}[colback=sky,colframe=navy,boxrule=0.9pt,arc=2mm,
    left=4mm,right=4mm,top=3mm,bottom=3mm]
  \textbf{Unit Overview.}\quad
  Unit 4 develops the mathematical language of chance.  Students begin by recognizing
  that patterns in data can arise from pure randomness, then build simulation skills
  to estimate probabilities empirically.  The probability rules half of the unit
  establishes sample spaces, complementary events, conditional probability,
  independence, and the addition rule --- the tools needed to calculate the likelihood
  of complex events.  The random-variables half introduces discrete probability
  distributions and their parameters (mean and standard deviation), rules for combining
  independent random variables, and two specific models --- the binomial and geometric
  distributions --- that describe counts of successes in repeated binary trials.
  This unit provides the probabilistic foundation for all of Units 5--9.
\end{tcolorbox}

\vspace{0.22in}

% ── Lessons in this unit ──────────────────────────────────────────────────────
\renewcommand{\arraystretch}{1.6}
\begin{skillbox}[Lessons in This Unit]{goldbox}
\begin{tabularx}{\linewidth}{@{} c >{\bfseries}p{0.46\linewidth} X @{}}
  \textbf{\#} & \textbf{Title} & \textbf{Focus} \\
  \hline
  4.1  & Introducing Statistics: Random and Non-Random Patterns?
       & Patterns in data vs.\ true randomness; questioning variation. \\
  \hline
  4.2  & Estimating Probabilities Using Simulation
       & Outcomes, events, simulation design; law of large numbers. \\
  \hline
  4.3  & Introduction to Probability
       & Sample spaces; classical probability; complement rule; long-run interpretation. \\
  \hline
  4.4  & Mutually Exclusive Events
       & Joint probability; disjoint events; Venn diagrams; $P(A \cap B) = 0$. \\
  \hline
  4.5  & Conditional Probability
       & $P(A\mid B)$; multiplication rule; two-way tables. \\
  \hline
  4.6  & Independent Events and Unions of Events
       & Independence vs.\ mutual exclusivity; addition rule $P(A \cup B)$. \\
  \hline
  4.7  & Introduction to Random Variables and Probability Distributions
       & Discrete random variables; probability tables and histograms; cumulative distributions. \\
  \hline
  4.8  & Mean and Standard Deviation of Random Variables
       & $\mu_X = \sum x\,P(X{=}x)$;\; $\sigma_X$; contextual interpretation. \\
  \hline
  4.9  & Combining Random Variables
       & Linear transformations; sums and differences of independent random variables. \\
  \hline
  4.10 & Introduction to the Binomial Distribution
       & BINS conditions; $P(X{=}k) = \binom{n}{k}p^k(1-p)^{n-k}$; binompdf. \\
  \hline
  4.11 & Parameters for a Binomial Distribution
       & $\mu_X = np$;\; $\sigma_X = \sqrt{np(1-p)}$;\; cumulative binomcdf; interpretation. \\
  \hline
  4.12 & The Geometric Distribution
       & Number of trials until first success; $P(X{=}k) = (1-p)^{k-1}p$;\; $\mu_X = 1/p$. \\
  \hline
\end{tabularx}
\end{skillbox}

\vspace{0.18in}

% ── Standards covered ─────────────────────────────────────────────────────────
\begin{skillbox}[AP Statistics Big Ideas \& Learning Objectives --- Unit 4]{greenbox}
\begin{tabularx}{\linewidth}{@{} >{\bfseries}p{0.14\linewidth} X @{}}
  VAR-1.F & Identify questions suggested by patterns in data; recognize that patterns
            do not necessarily imply non-random variation. \\
  UNC-2.A & Estimate probabilities using simulation; apply the law of large numbers. \\
  VAR-4.A & Calculate probabilities for events and their complements using sample spaces
            and the classical definition. \\
  VAR-4.B & Interpret probabilities as long-run relative frequencies. \\
  VAR-4.C & Explain why two events are (or are not) mutually exclusive using $P(A\cap B)$. \\
  VAR-4.D & Calculate conditional probabilities using $P(A\mid B) = P(A\cap B)/P(B)$
            and the multiplication rule. \\
  VAR-4.E & Calculate probabilities for independent events and for the union of two
            events using the addition rule. \\
  VAR-5.A & Represent the probability distribution for a discrete random variable
            as a table, graph, or function. \\
  VAR-5.B & Interpret a probability distribution in terms of shape, center, and spread. \\
  VAR-5.C & Determine mean and standard deviation of a discrete random variable. \\
  VAR-5.D & Interpret $\mu_X$ and $\sigma_X$ in context. \\
  VAR-5.E--F & Determine parameters for linear combinations and sums/differences
               of independent random variables. \\
  UNC-3.A & Determine whether a random variable follows a binomial distribution (BINS). \\
  UNC-3.B--C & Determine and interpret parameters of the binomial distribution. \\
  UNC-3.D--F & Calculate and interpret probabilities and parameters for the geometric
               distribution. \\
\end{tabularx}
\end{skillbox}

\end{document}
